% !TEX root = ../main.tex
\chapter{Connection fields (Zusammenhangsfelder)}\label{ch:c06}
\begin{paracol}{2}

\begin{sources}
\begin{tabularx}{\linewidth}{@{}l l L@{}}
\toprule
\textbf{Video} & \textbf{Length} & \textbf{Content}\\
\midrule
\seg{6}{1}{Lecture 6: Connection fields (Zusammenhangsfelder)} & 1:52:44 & Components as functions on the chart domain; the
naive directional derivative of a vector field and why it fails (extra term); linear connections (Zusammenh\"ange) and their
transformation law, flat regions; two ways to obtain a connection; the covariant derivative of a vector field, with the full
proof; the covariant derivative of an $(r,s)$ tensor field; affine manifolds; parallel vector fields, autoparallels, the
autoparallel equation and the acceleration.\\
\bottomrule
\end{tabularx}

\smallskip
\textbf{Transcript} \file{materials/transcripts/C06-V01.txt} (German YouTube auto-captions).
\end{sources}

\begin{reading}[Where this lecture sits in the other notes]
\emph{Geometric Anatomy}: lectures 21 to 25 (connections, covariant derivatives, parallel transport; lecture 25 was
already announced in \cref{ch:c05}). \emph{Gravity and Light} (TP-1 version): lectures 7 and 8 (the same material, there
with the metric and the Levi-Civita connection following).
\end{reading}

% ==================================================================
\section{Motivation: what is the acceleration?}\label{sec:cf-motivation}

After all the efforts of the previous lectures we take on further efforts, but today we also reap a fruit, the one promised at
the start of the course \ts{6}{1}{0:00}. Everybody knows from the second Newtonian axiom (``second Newtonian
action'' in the transcript) that mass times acceleration equals force; the force we shall model later by a vector field. But
we have already seen that the acceleration cannot be given in a chart by the second derivatives $\ddot q^a$ of the curve
components: these do not transform as vector components (\cref{ex:tf-second}), whereas the force does
\ts{6}{1}{0:42}.
At the end of this lecture we shall understand completely what the acceleration is and how to write it: one of the central
elements of Newtonian dynamics will be defined \ts{6}{1}{1:44}.
It is also essentially the end of the mathematical foundations: along the way we have collected everything needed in
classical mechanics, including general relativity \ts{6}{1}{2:04}.

The lecture is called \emph{connection fields}. They are not tensor fields (they transform differently), but with a
connection field and tensor fields one can do all sorts of things \ts{6}{1}{2:25}. The story starts with a
\emph{definition attempt}: we want a \textbf{directional derivative of a vector field}. We already have directional derivatives
of functions (a linear operator that eats a smooth function and gives a number); now we want ones acting on vector fields.
``The word attempt tells you already that it will go wrong; but it goes wrong in an interesting way, and it can be
repaired, by introducing connection fields'' \ts{6}{1}{3:08}.

% ==================================================================
\section{A change of convention for components}\label{sec:cf-convention}

A short recap: in \cref{ch:c05} the components of a vector field in a chart $(U,x)$ were component \emph{functions}
$X^i_{(x)}\colon x(U)\to\R$, depending on the point $x(p)$ of the chart image \ts{6}{1}{4:31}. All this is correct and
consistent, but while working out the exercises Schuller noticed that it leads in many places to needless ``chain rule
orgies'' \ts{6}{1}{5:54}. From now on the components are defined directly on the chart domain,
\[
  X^i_{(x)}\;:=\;\bigl(X^i_{(x)}\ \text{of lecture 5}\bigr)\circ x\ \colon\ U\to\R ,
\]
i.e.\ the old function after applying the chart map; the information is the same (one can always go back by composing with
$x^{-1}$), but this combination is the more convenient object \ts{6}{1}{6:37}. Everything of lecture~5 is unchanged except
that one no longer writes ``evaluated at $y(p)$'' on one side and ``evaluated at $x(p)$'' on the other: the transformation law
becomes an equation between functions on the overlap, with pointwise multiplication,
\[
  Y^a_{(y)}=\frac{\partial y^a}{\partial x^m}\,Y^m_{(x)} ,
\]
where $\partial y^a/\partial x^m$ now also denotes the function $\partial_m(y^a\circ x^{-1})\circ x$ on $U\cap V$.
``Philosophically nothing happens'' \ts{6}{1}{8:04}. Problem sheet~5 (handed out that day) already uses this convention;
exercise~1 on it shows the transformation in this form, and exercise~1b the product rule used below \ts{6}{1}{8:25}.

% ==================================================================
\section{The naive directional derivative of a vector field}\label{sec:cf-naive}

\subsection{The attempt}

\begin{definition}[New concepts: the naive directional derivative (``flower pot with flower'')]\label{def:cf-naive}
Let $(M,\mathcal O,\mathcal A)$ be a smooth manifold (this is always assumed unless said otherwise), $Y$ a vector field on $M$,
$X\in T_pM$ a tangent vector (not a field: one only needs the direction at the single point $p$) and $(U,x)\in\mathcal A$ a chart with
$p\in U$ \ts{6}{1}{9:47}. The \textbf{naive directional derivative} of $Y$ at $p$ in the direction $X$, written
$\nabla^{\text{naive}}_XY$, is defined by its components in this chart,
\begin{align*}
  \bigl(\nabla^{\text{naive}}_XY\bigr)^i_{(x)}(p)\;:=\;\lim_{\varepsilon\to0}\;
  \frac{Y^i_{(x)}\bigl(x^{-1}(x(p)+\varepsilon X_{(x)})\bigr)-Y^i_{(x)}(p)}{\varepsilon},
\end{align*}
where $X_{(x)}=(X^1_{(x)},\dots,X^d_{(x)})\in\R^d$ is the component tuple of $X$ in the chart \ts{6}{1}{11:10}.
\end{definition}
We use the symbol $\nabla$ (nabla) only as a label: ``forget everything you know about the nabla symbol''; it should differentiate
in the direction $X$ (the subscript) the object $Y$, and what one wants is again a vector field, so that the result has again
one index up \ts{6}{1}{11:31}.
The definition is made with respect to a chart (it could be made chart-independent, but this would be a ``technical overkill'').
Schuller decorates the symbol on the board with a flower: the pot \emph{with} the flower is the naive thing
(the flower marks ``naive''); later we shall give a definition where only the pot is left and the flower dropped
\ts{6}{1}{12:53}.

Where does the expression come from? Recall experimental physics. Go to the chart image \ts{6}{1}{13:40}. There the image
$x(p)$ of the base point and an arrow with the chart components $(X^1,\dots,X^d)$ can be drawn; and we are now interested,
\emph{in the chart}, in the point of $\R^d$ reached as $x(p)+\varepsilon X_{(x)}$ \ts{6}{1}{14:20}: we are on $\R^d$, so we may add,
namely a real number $\varepsilon$ times all components of $X$ \ts{6}{1}{15:01}. This is the chart image of a different point
of $M$, namely $q:=x^{-1}\bigl(x(p)+\varepsilon X_{(x)}\bigr)$ (back from the chart to the manifold with $x^{-1}$)
\ts{6}{1}{16:44}.
(A student corrects a slip in the index placement of a formula on the board; Schuller: ``test passed, thank you for the
correction'' \ts{6}{1}{16:02} [?].)

Next the vector field: in the chart $Y$ gives at every point a vector, with components $Y^i_{(x)}$ labelled by the point of $M$ it is
attached to; there is such a vector also at $q$ \ts{6}{1}{17:46}. How does the field change when we leave $p$ in the direction
$X$? ``There is no question how: of course like $\bigl(Y(q)-Y(p)\bigr)/\varepsilon$, with $\varepsilon\to0$'': that is
the difference quotient \ts{6}{1}{19:28}. Visually, however, we have subtracted a vector attached at $p$ from one attached at
$q$. ``How should I subtract from a vector attached there a vector attached here? Through years of indoctrination you have
learnt to accept that one simply \emph{shifts} it forward, by so-called parallel transport; nobody told you what it is, but it
sounds good.'' In the chart one can do it: one ignores the base point, takes the same components and subtracts them
\ts{6}{1}{20:53}. A smaller step $\varepsilon$ gives a smaller difference, and dividing by $\varepsilon$ one has the usual
difference quotient \ts{6}{1}{21:56}.

Written out, this is a familiar object: it is the $m$-th component of $X$ times the partial derivative in the $m$-th slot of the
$i$-th component of $Y$ in the chart (with exactly the decorations of the definition of the partial derivative):
\begin{equation}\label{eq:cf-naive}
  \bigl(\nabla^{\text{naive}}_XY\bigr)^i_{(x)}
  \;=\;X^m_{(x)}\,\frac{\partial}{\partial x^m}\,Y^i_{(x)}\;=\;X^m_{(x)}\,\partial_m\bigl(Y^i_{(x)}\circ x^{-1}\bigr)\circ x ,
\end{equation}
where $\partial/\partial x^m$ is the chart-induced basis vector acting on the function $Y^i_{(x)}\colon U\to\R$ \ts{6}{1}{22:41}.
So ``the flower pot with flower'' is exactly what was always said in experimental physics about the derivative of a vector field
in a direction \ts{6}{1}{24:26}.

\subsection{Why it fails}

The decisive question: we want the result to be a vector field. Are the $\bigl(\nabla^{\text{naive}}_XY\bigr)^i_{(x)}$ the
component functions of a vector field? ``I have said already that it goes wrong'' \ts{6}{1}{24:47}.

\begin{theorem}[The naive derivative is not a vector field]\label{thm:cf-naive-fail}
Under a change to a chart $(V,y)$, the components \eqref{eq:cf-naive} transform as
\begin{align*}
  \bigl(\nabla^{\text{naive}}_XY\bigr)^a_{(y)}
  &=\frac{\partial y^a}{\partial x^q}\,\bigl(\nabla^{\text{naive}}_XY\bigr)^q_{(x)}\
  &\quad+\frac{\partial^2 y^a}{\partial x^m\partial x^q}\,X^m_{(x)}\,Y^q_{(x)} ,
\end{align*}
so they are the components of a vector field only if the second term vanishes, which in general it does not.
\end{theorem}
\begin{proof}
To check whether the components of a quantity form a vector field one writes the same definition in an arbitrary other chart and
reduces it to the components in chart $x$: if the vector law comes out, they are components of a vector field
\ts{6}{1}{25:52}. So compute, by definition in the chart $y$,
$(\nabla^{\text{naive}}_XY)^a_{(y)}=X^m_{(y)}\,\dfrac{\partial}{\partial y^m}\,Y^a_{(y)}$ \ts{6}{1}{26:54}. Two known facts are inserted:
the vector $X$ expanded in the two bases, $X^m_{(y)}\dfrac{\partial}{\partial y^m}=X^k_{(x)}\dfrac{\partial}{\partial x^k}$ (the basis
transformation of \cref{ch:c04}; in the lecture it is written the other way round, $\partial/\partial y^m=(\partial x^l/\partial y^m)\,\partial/\partial x^l$, and
combined with the transformation law of $X^m$) \ts{6}{1}{27:35}, and the transformation law $Y^a_{(y)}=(\partial y^a/\partial x^q)\,Y^q_{(x)}$. The
Kronecker delta $\delta^k{}_l=\dfrac{\partial y^m}{\partial x^l}\dfrac{\partial x^k}{\partial y^m}$ of \cref{thm:tf-inverse} reduces the sums
(``one always asks: is there a simplification? the trained eye sees a repeated index up and down'') \ts{6}{1}{30:03}. This gives
\[
  \bigl(\nabla^{\text{naive}}_XY\bigr)^a_{(y)}=X^l_{(x)}\frac{\partial}{\partial x^l}\Bigl(\frac{\partial y^a}{\partial x^q}\,Y^q_{(x)}\Bigr).
\]
Now $\partial/\partial x^l$ acts on a product of two functions: the product rule (problem sheet~5, exercise~1b) gives two terms
\ts{6}{1}{32:06}. In the first, it acts on $Y^q_{(x)}$:
\[
  \frac{\partial y^a}{\partial x^q}\;X^l_{(x)}\frac{\partial}{\partial x^l}Y^q_{(x)}
  =\frac{\partial y^a}{\partial x^q}\bigl(\nabla^{\text{naive}}_XY\bigr)^q_{(x)} ,
\]
which has exactly the form of the vector law. In the second, it acts on the Jacobian: one must watch the order, the derivative
$\partial y^a/\partial x^q$ was taken first and then $\partial/\partial x^l$ is applied from the left, so
\[
  X^l_{(x)}\,Y^q_{(x)}\;\frac{\partial^2y^a}{\partial x^l\,\partial x^q} .
\]
\ts{6}{1}{33:30}. This second term is an \emph{extra term}. ``If it were not there, the object would be a vector field; but it is
there'' \ts{6}{1}{34:54}.
\end{proof}
So the extra term destroys the tensor property: ``this is actually the end of the proof: no, these are not the component functions
of a vector field'' \ts{6}{1}{35:36}.
It involves second derivatives of the transition map; compare with the extra term of $\ddot q^i$ in \cref{ex:tf-second}.
\begin{coursenote}[Why $\nabla^{\text{naive}}$ looks so natural, and what is missing]
The extra term is linear in $X$ and in $Y$ and built only from the second derivative of the transition map: it vanishes for
linear changes of chart (e.g.\ between Cartesian systems) but not for curvilinear ones. That is why, in
Cartesian coordinates, the old recipe of experimental physics ``works''. (Added; not said in class.)
\end{coursenote}
How can one do better? ``There is no way to define this better, except, and this is the point, by having the courage to introduce
\emph{new structure}'': on a bare smooth manifold you cannot. One defines a new structure (and then has to \emph{choose} it:
you know the game by now), and it goes by the name that gave the lecture its title \ts{6}{1}{35:56}.

% ==================================================================
\section{Linear connections}\label{sec:cf-connections}

The new structure is a \textbf{linear connection} (linearer Zusammenhang). Of all tensor fields we need only connections with
one index up and two down, ``and this suffices forever'' \ts{6}{1}{36:39}. Indices refer to a chart, so the definition is made
with respect to a chart \ts{6}{1}{37:21}.

\begin{definition}[New concepts: linear connection, connection coefficients]\label{def:cf-connection}
Let $(M,\mathcal O,\mathcal A)$ be a smooth manifold of dimension $d$. A \textbf{linear connection} (also: connection field,
\emph{Zusammenhang}) on $M$ is, for every chart $(U,x)$ of the atlas, a set of $d^3$ smooth functions
\[
  \Gamma^a{}_{bc}{}_{(x)}\colon U\to\R\qquad(a,b,c=1,\dots,d)
\]
(the \textbf{connection coefficient functions}, or simply ``the $\Gamma$'s'', the \emph{Christoffel symbols} when taken from a metric;
one index up, two down), such that under a change to a chart $(V,y)$ they transform, on the overlap $U\cap V$, as
\begin{align}\label{eq:cf-gamma-trafo}
  \Gamma^a{}_{bc}{}_{(y)}
  &=\frac{\partial y^a}{\partial x^m}\,\frac{\partial x^n}{\partial y^b}\,\frac{\partial x^p}{\partial y^c}\;\Gamma^m{}_{np}{}_{(x)}\notag\\
  &\quad+\;\frac{\partial y^a}{\partial x^k}\;\frac{\partial^2x^k}{\partial y^b\,\partial y^c}
\end{align}
\ts{6}{1}{38:07}. The first term is the transformation law of a $(1,2)$ tensor field; the \textbf{inhomogeneous} second term
depends only on the transition map, not on $\Gamma$.
\end{definition}
(The sentence ``$d^3$ functions'' is in the transcript as ``d to the third''.)
They are \emph{not} the components of a $(1,2)$ tensor field: were the second term absent they would be, and no new name would be
needed. The whole point is that this term is there \ts{6}{1}{39:33}. The law holds at all points of the overlap of the two charts;
nothing can be said about points outside it \ts{6}{1}{42:09}.
The $\Gamma$ appear on the right-hand side in the first term, but the second term depends only on the transition functions
\ts{6}{1}{42:56}.

\begin{intuition}[What the inhomogeneous term means: flat regions, and why Cartesian is not enough]
A wise guy might say: ``I define $\Gamma$ in one coordinate system to be zero everywhere.'' If this is possible on a whole region
of the manifold, in a suitable chart, the region is called \textbf{flat} (flach). In Newtonian physics space is flat, so why
do we need all this differential geometry; we just put $\Gamma=0$? Fine; then go to another chart $y$. The tensor part stays $0$, but
the extra term appears as soon as the coordinate change is \emph{not linear}: e.g.\ from Cartesian to polar coordinates the
second derivative does not vanish. So you cannot demand $\Gamma=0$ in all coordinate systems, even in flat space, and why
should we deprive ourselves of polar coordinates, which one uses all the time for rotation-symmetric problems? Then the $\Gamma$'s
change even if they were all zero before \ts{6}{1}{43:17}.
For a tensor this is not so: if the components of a tensor vanish in one chart they vanish in every chart, since a transformation
law that is linear in the components maps $0$ to $0$. So a connection is ``a truly interestingly transforming object''
\ts{6}{1}{44:40}. These remarks become clearer in this and the next lecture \ts{6}{1}{45:21}.
\end{intuition}
\begin{coursenote}[On the word ``flat'']
In the lecture ``flat region'' is explained only through the choice $\Gamma=0$ in a suitable chart, as above. The curvature
(and torsion) of a connection are not defined in this lecture; the notion of flatness is therefore used here only in
this loose sense. (Added; not said in class.)
\end{coursenote}

\subsection{How to obtain a connection}\label{ssec:cf-obtain}

The definition looks very artificial: ``a collection of components, and if you go to a new chart you do this and that''. But one
can get such objects \ts{6}{1}{45:41}. There are two ways, as we had already for tensor fields and vector fields
\ts{6}{1}{46:02}.
\begin{enumerate}
  \item \textbf{Prescribe} them: in every chart of the atlas $\mathcal A$ choose $d^3$ smooth functions (they only need to be
        smooth); the only thing to ensure is that on the overlaps of two charts law \eqref{eq:cf-gamma-trafo} holds, otherwise the definition would
        be contradictory between the charts \ts{6}{1}{46:22}. One says one \emph{chooses} a connection. If one can choose all
        $\Gamma=0$ in a whole chart region, the region is flat, as above; if this is not possible (``by far the more frequent
        case'') the manifold is not flat; we come back to this \ts{6}{1}{47:46}.
  \item \textbf{Construct} them: from tensor fields. ``Believe it or not, one can even define them from tensor fields; one of the
        most famous methods is from a \emph{symmetric $(0,2)$ tensor field}'' (this is problem sheet~5, exercise~2; one must then
        check that the behaviour \eqref{eq:cf-gamma-trafo} really holds) \ts{6}{1}{48:27}.
\end{enumerate}
\begin{aside}[Why $\Gamma$ matters: the link with the naive derivative]
Why all this? The $\Gamma$'s do not transform as tensors, in a way that lets one add a suitable term to repair the directional
derivative of a vector field: ``with such objects, which transform completely differently, one can make that directional
derivative work, by adding a corresponding term'' \ts{6}{1}{49:08}.
\end{aside}
A natural worry: for a vector field it works, but what about the directional derivative of a $(0,1)$ or general $(r,s)$ field, does one
need a new object each time? ``The good news: no. Even the directional derivative of a \emph{general} $(r,s)$ tensor field is
completely definable with the \emph{same} linear connection coefficients, one index up and two down; they are merely used in a
convoluted way. So one needs only once to choose such a structure and one has the derivative of every tensor field under control''
\ts{6}{1}{50:11}.
Do not confuse: in lecture~4 we already had the directional derivative \emph{acting on functions} (the $(0,0)$ case) and the vectors
themselves, a $(1,0)$ field; as soon as one acts on a tensor field of higher rank which is not just $(0,0)$ one needs the additional
structure. The directional derivative of a function will then be redefined accordingly \ts{6}{1}{50:52}.
After a break \ts{6}{1}{51:12}.

% ==================================================================
\section{The covariant derivative}\label{sec:cf-covariant}

Now ``we reap the fruits''. First of a vector field, then of general tensor fields \ts{6}{1}{51:33}.

\begin{definition}[New concepts: covariant derivative of a vector field]\label{def:cf-covariant}
Let $\Gamma$ be a linear connection on $M$. The \textbf{covariant derivative} (kovariante Ableitung, here: the ``ordinary'' directional
derivative, with the flower dropped) $\nabla_XY$ of a vector field $Y$ at $p\in M$ in the direction $X\in T_pM$ is defined,
with respect to a chart $(U,x)$ with $p\in U$, by its components
\begin{align}\label{eq:cf-covariant}
  \bigl(\nabla_XY\bigr)^i_{(x)}&:=\bigl(\nabla^{\text{naive}}_XY\bigr)^i_{(x)}+\Gamma^i{}_{jm}{}_{(x)}\,Y^j_{(x)}\,X^m_{(x)}\notag\\
  &\;=X^m\frac{\partial Y^i}{\partial x^m}+\Gamma^i{}_{jm}\,Y^j\,X^m
\end{align}
(all components with respect to $x$ and at the point $p$; letting $p$ run, one gets for a vector field $X$ again a vector field)
\ts{6}{1}{52:18}. Mind the \textbf{order of the indices}: $X$, the direction in which one differentiates, goes with the
\emph{last} index of $\Gamma$; $Y$, the field being differentiated, goes with the \emph{second to last} index \ts{6}{1}{54:06}.
\end{definition}
The result is a vector field; and the definition is always \emph{with respect to a chosen linear connection}: with another
connection the derivative looks different \ts{6}{1}{56:31}. ``The $\Gamma$'s are secret information: they tell you how, in a given
chart, the vector components change under that parallel transport which we carried out so naively, i.e.\ they are the
coefficients/components of a parallel transport'' \ts{6}{1}{57:13}. The ``flower pot without flower'' is very useful (one can plant
flowers in it); from now on one writes $\nabla_XY$ without the flower, ``we are not that silly'' \ts{6}{1}{55:09}.

\begin{theorem}[$\nabla_XY$ is a vector field]\label{thm:cf-covariant}
The components \eqref{eq:cf-covariant} transform with the vector law under chart changes: $\nabla_XY$ is a vector field.
\end{theorem}
\begin{proof}
The announced horrendous computation: ``it helps nothing'' \ts{6}{1}{54:47}. We must be careful and concentrated
\ts{6}{1}{55:30}. In the chart $y$ the definition reads
$(\nabla_XY)^a_{(y)}=(\nabla^{\text{naive}}_XY)^a_{(y)}+\Gamma^a{}_{bc}{}_{(y)}Y^b_{(y)}X^c_{(y)}$
(for the connection chosen once) \ts{6}{1}{57:55}, and we replace every $y$-quantity by the corresponding $x$-quantity, choosing indices
cleverly to avoid collisions \ts{6}{1}{58:15}.
By \cref{thm:cf-naive-fail} the first term is
\[
  \bigl(\nabla^{\text{naive}}_XY\bigr)^a_{(y)}
  =\frac{\partial y^a}{\partial x^q}\bigl(\nabla^{\text{naive}}_XY\bigr)^q_{(x)}
  +\frac{\partial^2y^a}{\partial x^l\partial x^q}\,X^l_{(x)}Y^q_{(x)}.
\]
For the second term we insert the transformation law \eqref{eq:cf-gamma-trafo} of the connection and those of the vector
components,
\begin{align*}
  \Gamma^a{}_{bc}{}_{(y)}Y^b_{(y)}X^c_{(y)}
  &=\Bigl(\frac{\partial y^a}{\partial x^i}\frac{\partial x^j}{\partial y^b}\frac{\partial x^k}{\partial y^c}\,\Gamma^i{}_{jk}{}_{(x)}
   +\frac{\partial y^a}{\partial x^q}\frac{\partial^2x^q}{\partial y^b\partial y^c}\Bigr)\\
  &\qquad\cdot\frac{\partial y^b}{\partial x^r}Y^r_{(x)}\,
   \frac{\partial y^c}{\partial x^s}X^s_{(x)}\\
  &=\frac{\partial y^a}{\partial x^i}\,\Gamma^i{}_{rs}{}_{(x)}\,Y^r_{(x)}X^s_{(x)}\\
  &\quad+\frac{\partial y^a}{\partial x^q}\frac{\partial^2x^q}{\partial y^b\partial y^c}\frac{\partial y^b}{\partial x^r}\frac{\partial y^c}{\partial x^s}\,Y^r_{(x)}X^s_{(x)} .
\end{align*}
In the first term of the second line, the products $\dfrac{\partial x^j}{\partial y^b}\dfrac{\partial y^b}{\partial x^r}=\delta^j{}_r$ and the
analogous one for $c$ were used: ``an orgy of Kronecker deltas'', the boxes of indices can be forgotten once they are replaced by deltas
\ts{6}{1}{1:05:01}. (The lecture has index collisions at this point, e.g.\ the summation index $k$ was used twice;
Schuller renames them on the board: ``this happens to everybody'' \ts{6}{1}{1:07:47}.)
Collecting all terms,
\begin{align*}
  (\nabla_XY)^a_{(y)}
  &=\frac{\partial y^a}{\partial x^q}\Bigl[\bigl(\nabla^{\text{naive}}_XY\bigr)^q_{(x)}+\Gamma^q{}_{rs}{}_{(x)}Y^r_{(x)}X^s_{(x)}\Bigr]\
  &\quad+\Bigl[\frac{\partial^2y^a}{\partial x^s\partial x^r}\\
  &\qquad+\frac{\partial y^a}{\partial x^q}\frac{\partial^2x^q}{\partial y^b\partial y^c}\frac{\partial y^b}{\partial x^r}\frac{\partial y^c}{\partial x^s}\Bigr]Y^r_{(x)}X^s_{(x)} .
\end{align*}
The first line is exactly $\partial y^a/\partial x^q\,(\nabla_XY)^q_{(x)}$, the vector law \ts{6}{1}{1:08:49}; if the second line were
absent we would be done. It is zero, as the following side computation shows \ts{6}{1}{1:09:31}.

\emph{Side computation} (the two second derivatives have a $+$ between them, one would like a $-$). Start from what is known, the
inverse Jacobians (\cref{thm:tf-inverse}):
\[
  \frac{\partial x^q}{\partial y^b}\,\frac{\partial y^b}{\partial x^r}=\delta^q{}_r .
\]
It holds in every chart (the right-hand side consists of numbers), so we may act with $\partial/\partial x^s$ on both sides; the
product rule gives
\[
  \frac{\partial}{\partial x^s}\Bigl(\frac{\partial x^q}{\partial y^b}\Bigr)\frac{\partial y^b}{\partial x^r}
  +\frac{\partial x^q}{\partial y^b}\,\frac{\partial^2y^b}{\partial x^s\partial x^r}=0
\]
(mind the order: $\partial y^b/\partial x^r$ was differentiated first, $\partial/\partial x^s$ then comes from the left)
\ts{6}{1}{1:10:32}.
In the first term the derivative is with respect to $x^s$ but of a function of the form $\partial x^q/\partial y^b$, which one
would like to differentiate with respect to $y$: express $\partial/\partial x^s=(\partial y^c/\partial x^s)\,\partial/\partial y^c$, so that
\[
  \frac{\partial y^c}{\partial x^s}\frac{\partial^2x^q}{\partial y^c\partial y^b}\frac{\partial y^b}{\partial x^r}
  +\frac{\partial x^q}{\partial y^b}\frac{\partial^2y^b}{\partial x^s\partial x^r}=0 .
\]
Multiply by $\partial y^a/\partial x^q$ and sum over $q$: the second term becomes $\delta^a{}_b\,\partial^2y^b/\partial x^s\partial x^r$, so
the identity is exactly $R=0$ \ts{6}{1}{1:15:47}. (In the lecture the steps are done with $\partial/\partial x^m$ and a multiplication by
$\partial y^b/\partial x^v$; one ends in the same identity. ``Replace the indices, check there is no collision, and the whole second line
disappears'' \ts{6}{1}{1:16:51}.)

\begin{coursenote}[Schwarz's theorem only within one chart]
Schuller warns for all time: the theorem of Schwarz (partial derivatives of a smooth function commute) inherits to these expressions
\emph{only for coordinates of the same chart}: $\partial^2x^q/\partial y^b\partial y^c$ is symmetric in $b,c$, but mixed expressions with
$y$ and $x$ derivatives (e.g.\ $\partial^2y/\partial x\partial x$ against $\partial^2x/\partial y\partial y$) cannot be swapped freely
\ts{6}{1}{1:12:58}. Also, the derivatives here are not partial derivatives of functions on $\R^d$ but chart derivatives
(\cref{ch:c05}).
\end{coursenote}

As the second line vanishes, only the vector law remains: $\nabla_XY$ is indeed a vector field \ts{6}{1}{1:17:34}.
\end{proof}
``If you really follow every step, which will cost you a whole afternoon, you have learnt a lot'': this is your ``journeyman's
examination'' in differential geometry \ts{6}{1}{1:17:13}. Are there questions? None \ts{6}{1}{1:17:54}.

\begin{aside}[Why do this by hand]
As a motivation: for the rest of your career, in at least three further theory lectures, you will relearn the same arguments, without
recognising that they are the same. The divergence of a vector field in spherical polar coordinates? One looks it up in a formula
collection, but it is just computing the $\Gamma$'s in polar coordinates; the curl, all this ``rubbish of electrodynamics'' is the same
thing. ``This is the minimal story, reduced to the essential core; general relativity is said to be mathematically hard, but this
is it; if you understand it you have understood a lot of the technique of theoretical physics.'' That is why it is done once here
\ts{6}{1}{1:18:32}. And a remark on the exam: no aids are allowed, ``I mean it''; everything done here is so important that you must
be able to do it by heart; after half a year it flows, start today \ts{6}{1}{1:02:34}.
\end{aside}

\subsection{Covariant derivative of an $(r,s)$ tensor field}\label{ssec:cf-tensor}

The vector field was only a special case (but one that demanded all the cruelty); for general tensor fields it is just as easy
\ts{6}{1}{1:20:17}.

\begin{definition}[New concepts: covariant derivative of an $(r,s)$ tensor field]\label{def:cf-covariant-tensor}
Let $T$ be an $(r,s)$ tensor field and $X$ a \emph{vector}: one can ask how a tensor field changes in a direction, but the direction
remains a vector, ``not a tensor'' \ts{6}{1}{1:21:24}. The covariant derivative $\nabla_XT$ is again an $(r,s)$ tensor field,
defined in a chart $(U,x)$ by its components
\begin{align*}
  \bigl(\nabla_XT\bigr)^{i_1\dots i_r}{}_{j_1\dots j_s}
  &=X^m\,\frac{\partial}{\partial x^m}\,T^{i_1\dots i_r}{}_{j_1\dots j_s}\\
  &\quad+\sum_{k=1}^r\Gamma^{i_k}{}_{t\,m}\,T^{i_1\dots t\dots i_r}{}_{j_1\dots j_s}\,X^m
        \;-\;\sum_{l=1}^s\Gamma^{t}{}_{j_l\,m}\,T^{i_1\dots i_r}{}_{j_1\dots t\dots j_s}\,X^m ,
\end{align*}
all components with respect to $x$. The first line is the naive part, acting on the bare component functions \ts{6}{1}{1:22:46}.
The rule: for \emph{each upper index} $i_k$ of $T$ there is an extra term with a plus sign, in which that index is replaced by a dummy
$t$ (which then appears in the $k$-th upper slot of $T$) and $\Gamma^{i_k}{}_{t\,m}$ has the replaced index as its \emph{upper} index; for \emph{each lower index}
$j_l$ there is an extra term with a \emph{minus} sign, the index $j_l$ becoming the first lower index of
$\Gamma^{t}{}_{j_l\,m}$ and $t$ sitting in the $l$-th lower slot of $T$ \ts{6}{1}{1:23:27}. In all cases the last index of $\Gamma$ is
reserved for the vector $X$ in whose direction one differentiates and is summed with $X^m$ \ts{6}{1}{1:26:55}.
\end{definition}
(Schuller's naming of the dummy indices on the board, ``$t$ in a different system'' etc., was done aloud with some remarks that the
index of an index should not be confused with the index itself: ``only a madman would mix an index with the index of an index''
\ts{6}{1}{1:24:30}.)
So in short: every upper index behaves like a vector index (plus), every lower index like a covector index (minus). ``This you must
know in your sleep; then it is simple'' \ts{6}{1}{1:28:17}. Why plus and minus? ``One could derive it very nicely if the whole abstract
theory on the manifold had been done; but honestly, in dirty practice, in exams, this is what we need'' \ts{6}{1}{1:28:38}.
Now you know how to covariantly differentiate any tensor field; the important proviso is that this always happens
\emph{with respect to a chosen linear connection} \ts{6}{1}{1:29:19}.

\paragraph{A student asks for a numerical example.} A student asks for a concrete example with numbers (a matrix, say). Answer: all this is
theory and it is correct; as in school one can compute derivatives of functions in a concrete case, and Schuller makes a point
of \emph{not} giving one now. What does the ``times'' between $\Gamma$ and $T$ mean? It is not the product on the reals, it is not a
tensor product: it is the \textbf{pointwise product of two functions} $U\to\R$ (the $\Gamma^{a}{}_{bc}$ being $d^3$ component functions
of the connection, the $T$'s component functions of a tensor field); at a point it is the product of the evaluated factors
\ts{6}{1}{1:31:47}. So an example in class would implicitly cover many things that become explicit only if you do it yourself.
``You have to work at it at home until you find out how to compute it; only in this process of finding out do you learn mathematics,
and only in it do you learn physics. It is natural at university that you get the theory (and I give examples now and then),
but the problem sheets \emph{are} the examples; you must fall over every stumbling block once. You must try to crack every
single exercise yourself, ten hours a week if need be, others manage in two; you can write a letter to the minister of
education. Afterwards you are a really good physicist, world class, since there is nothing more to know'' \ts{6}{1}{1:33:51}.
``You have to train'' \ts{6}{1}{1:35:34}.

% ==================================================================
\section{Affine manifolds}\label{sec:cf-affine}

\begin{definition}[New concepts: affine manifold]\label{def:cf-affine}
An \textbf{affine manifold} (affine Mannigfaltigkeit) is a quadruple $(M,\mathcal O,\mathcal A,\nabla)$: a set $M$ with topology
$\mathcal O$, a smooth atlas $\mathcal A$, and a covariant derivative $\nabla$ (equivalently, a linear connection, i.e.\ in
every chart of the atlas $d^3$ functions $\Gamma^a{}_{bc}$, compatible on overlaps) \ts{6}{1}{1:35:54}.
\end{definition}
One has to \emph{put on the table} a huge amount of information, the $d^3$ functions in every chart; someone else may come with
other functions on every chart (compatible on overlaps) and that is his covariant derivative: the two then have different affine
manifolds, so the choice must be made \ts{6}{1}{1:37:29}. ``And which choice should I make?'' Physics says: for instance the first
Newtonian axiom. As said in the first lecture, the point of the first axiom is not how particles move if no force acts, but that
particles under certain circumstances move such that one may say no force acts on them, and these motions are straight lines,
``and, as we shall show later, what that has to do with the covariant derivative''. So physics, or experiment, tells us how the
$\Gamma$'s have to be chosen \ts{6}{1}{1:38:31}.
Further material is on problem sheet~5: exercises 2, 3, 4 (``from now on it is only terminology'') \ts{6}{1}{1:39:33}.

% ==================================================================
\section{Parallel transport and autoparallels}\label{sec:cf-parallel}

In the last twelve minutes we redeem the promise made at the beginning \ts{6}{1}{1:40:24}.

\begin{definition}[New concepts: parallel vector field, autoparallel vector field]\label{def:cf-parallel}
Let $(M,\mathcal O,\mathcal A,\nabla)$ be an affine manifold, $X$ a vector field.
\begin{enumerate}
  \item A vector field $Y$ is called \textbf{parallel along} (``parallel transported along'', as the literature says, although
        nothing is transported, ``one thinks something happens, it is simply parallel'') the vector field $X$ if
        $\nabla_XY=0$ at all points of $M$ \ts{6}{1}{1:41:05}.
  \item A vector field $X$ is called \textbf{autoparallel} (parallel to itself) if $\nabla_XX=0$ \ts{6}{1}{1:42:30}.
\end{enumerate}
\end{definition}
The intuition for the first: the change of $Y$ in the direction of $X$ is zero, wherever you are; $Y$ does not change along $X$
\ts{6}{1}{1:41:48}. Whether $Y$ is parallel along $X$ depends on the chosen linear connection \ts{6}{1}{1:42:09}.
For the second: ``follow your nose'', as the children's joke goes. At every point where I stand my nose defines a vector (a vector
field); if I walk in the direction of my nose I walk along a straight line. A straight line is often defined as the shortest
connection, but this needs a metric (a measure of length) \ts{6}{1}{1:43:12}. With a connection alone (equivalent to a covariant
derivative) one cannot say what the \emph{shortest} lines are (we shall do this later), but one can say what the \emph{straight lines} are:
those whose tangent vector field is autoparallel \ts{6}{1}{1:43:53}.

In components, with respect to a chart $x$: take $X^m\,\partial/\partial x^m$, let it act on the components of $X$ and add the correction,
\begin{equation}\label{eq:cf-autoparallel-field}
  X^m_{(x)}\,\frac{\partial}{\partial x^m}X^a_{(x)}\;+\;\Gamma^a{}_{bc}{}_{(x)}\,X^b_{(x)}X^c_{(x)}=0
\end{equation}
(simply the covariant derivative of the vector field $Y=X$ in the direction $X$) \ts{6}{1}{1:44:36}.

With vector fields one cannot think so well; we are used to curves. So the same for a curve:
\begin{definition}[Autoparallel curve]\label{def:cf-autoparallel-curve}
A curve $\gamma\colon I\to M$ on an affine manifold is called an \textbf{autoparallel} if its tangent vector (at every point of the
curve a new vector, i.e.\ a vector field \emph{along the curve}) is autoparallel \ts{6}{1}{1:45:17}.
\end{definition}
(The curve gives a vector field only along its image, not on all of $M$; the defining condition then asks the covariant derivative of the
tangent in the direction of the tangent itself; this makes sense because the derivative only needs the direction at the point and the
values of the field along the curve. Added, not said in class.)

Write it in a chart \ts{6}{1}{1:45:59}. Let the curve pass through the chart domain: the part of the parameter interval mapped into $U$ is
the preimage; $x\circ\gamma$ is the chart representation, with components $q^a(\lambda)=x^a(\gamma(\lambda))$
\ts{6}{1}{1:46:19}. The tangent vector has, as known from experimental physics and from \cref{ch:c04}, the components
$\dot q^a(\lambda)=(q^a)'(\lambda)$ \ts{6}{1}{1:47:01}. Substituting these as the components $X^a$ in \eqref{eq:cf-autoparallel-field},
using the chain rule backwards to turn $X^m\partial_m(\dot q^a)$ into $\ddot q^a$ (``Newton would be proud of us''),
\ts{6}{1}{1:47:42}:
\begin{theorem}[Autoparallel equation]\label{thm:cf-autoparallel}
A curve is autoparallel if and only if, in every chart, its component functions satisfy
\[
  \ddot q^a(\lambda)\;+\;\Gamma^a{}_{bc}\bigl(q(\lambda)\bigr)\,\dot q^b(\lambda)\,\dot q^c(\lambda)\;=\;0 ,
\]
with the $\Gamma^a{}_{bc}$ evaluated at the point of the curve at parameter $\lambda$ (``all this, let me leave out in the formula'')
\ts{6}{1}{1:48:24}.
\end{theorem}
(In the lecture the equation is written $\dfrac{d}{d\lambda}\dot q^a+\Gamma^a{}_{bc}\,\dot q^b\dot q^c=0$, ``if you are very
physics-affine write dots instead of derivatives'', with the parameter $\lambda$ possibly thought of as time \ts{6}{1}{1:49:09}.)

\paragraph{The acceleration.} If you massage the first term, by the chain rule backwards it is nothing but the second derivative of the chart
representation of the curve: ``we know that this alone has no vector components'' (as repeatedly said, on the problem sheet too); but
\emph{with} the additional term it has vector components. That is the point: only the whole left-hand side, the second derivative \emph{plus}
the correction quadratic in the velocity, defines the \textbf{components of the acceleration} \ts{6}{1}{1:49:50}.
\begin{equation*}
  a^a\;:=\;\ddot q^a+\Gamma^a{}_{bc}\,\dot q^b\dot q^c
\end{equation*}
(the notation $a^a$ is added; not said in class). This is the acceleration of Newton's theory, the missing piece of $m\,a=F$.
What then does ``the acceleration is zero'' mean? It means that the curve is an autoparallel, a straight line. Newton says in the second
axiom: the \emph{deviation from straight motion} is proportional to the force; with no force there is no deviation, and the motion is
the straight one \ts{6}{1}{1:50:53}.
But beware: it is the straight line in the sense of ``following one's nose'', \emph{not} the straight line as the shortest
connection (formally one can relate the two, but this is subtle; only with additional structure) \ts{6}{1}{1:51:13}.

\begin{secondary}[Outlook given at the end]
This is the end of the mathematical foundations. Schuller reserves the right to do, on Thursday, a big worked example ``with all
the trimmings'', or otherwise to start directly using all the mathematics to write the physical notions (spacetime, time, absolute time)
in mathematical terms and then to formulate the Newtonian axioms correctly \ts{6}{1}{1:51:56}.
\end{secondary}

% ==================================================================
\section*{Key concepts}
\begin{keyconcepts}
\textbf{Retrieval questions.}
\begin{enumerate}
  \item How is the naive directional derivative of a vector field defined, and why is it not a vector field? What is the extra term?
  \item State the transformation law of a linear connection. Which part is tensorial and which not? Why can $\Gamma=0$ not hold in all charts?
  \item Which two ways are there to obtain a connection?
  \item Define the covariant derivative of a vector field; which index of $\Gamma$ goes with the direction, which with the field?
  \item Sketch the proof that $\nabla_XY$ is a vector field (which identity makes the second line vanish?).
  \item Write the covariant derivative of an $(r,s)$ tensor field: sign and index position for upper and lower indices.
  \item What is an affine manifold? Which choices does it contain, and who makes them?
  \item Define parallel and autoparallel; write the autoparallel equation and say what the acceleration is.
\end{enumerate}
\textbf{Concepts.}
\begin{itemize}
  \item Naive derivative and its failure (\cref{def:cf-naive,thm:cf-naive-fail}).
  \item Linear connection and its law (\cref{def:cf-connection}); flat regions; ways to obtain one (\cref{ssec:cf-obtain}).
  \item Covariant derivative (\cref{def:cf-covariant,thm:cf-covariant,def:cf-covariant-tensor}); affine manifold (\cref{def:cf-affine}).
  \item Parallel, autoparallel, autoparallel equation, acceleration (\cref{def:cf-parallel,thm:cf-autoparallel}).
\end{itemize}
\end{keyconcepts}

\flushside
\end{paracol}
